% !TeX root = ../main.tex
\section{Methodology}\label{sec:method}

\subsection{Platform}
\newcommand{\acspinning}{single-thread runs and the loaded-latency probe pinned to CPU~2 (a physical
core, not CPU~0); multi-threaded runs place one thread per physical core}
% Memory is covered by the shared table's single Memory row (MemTotal; DIMMs from a saved dmidecode).
All measurements were taken on the ten-core i9-10900F desktop (Comet Lake, Skylake-family core, one NUMA
node) in \cref{tab:environment}, generated by \code{scripts/collect_env.sh}; each raw-data directory also
stores an \code{env.txt} snapshot taken just before that experiment.
%% Measurement environment table (common requirement: measurement discipline).
%% Values come from common/system_info.tex. Before \input-ing this file a project may define
%%   \acspinning   how runs were pinned/isolated
%%   \acsenvextra  extra rows, e.g.  \newcommand{\acsenvextra}{SSD & model, firmware, PCIe link \\}
\providecommand{\acspinning}{\todo{e.g.\ \code{taskset -c 2} on a P-core; isolcpus if used}}
\providecommand{\acsenvextra}{}
\begin{table}[tbp]
  \centering\small
  \caption{Measurement environment.}\label{tab:environment}
  \begin{tabularx}{\linewidth}{@{}>{\raggedright\arraybackslash}p{0.2\linewidth}X@{}}
    \toprule
    CPU               & \sysCPU{} (\sysArch) \\
    SIMD ISA          & \sysISA \\
    Topology          & \sysCores; NUMA: \sysNUMA \\
    SMT               & \sysSMT \\
    Caches            & \sysCaches \\
    Memory            & \sysMemory; DIMMs: \sysDIMM \\
    Pages             & base \sysPageSize; THP: \sysHugePages \\
    OS / kernel       & \sysOS; \sysKernel \\
    Frequency policy  & governor: \sysGovernor; driver: \sysDriver; turbo: \sysTurbo; EPP: \sysEPP \\
    Pinning / isolation & \acspinning; \sysIsolation \\
    Compilers         & \sysGCC; \sysClang \\
    Tools             & perf: \sysPerf{} (perf\_event\_paranoid \sysPerfParanoid); FIO: \sysFio; MLC: \sysMLC \\
    GPU (optional)    & \sysGPU; CUDA \sysCUDA \\
    \acsenvextra
    Env.\ snapshot     & \sysCollected{} (each raw run directory also stores its own \code{env.txt}) \\
    \bottomrule
  \end{tabularx}
\end{table}


\subsection{Tools}
All benchmarks share the harness in \code{common/code} (raw monotonic timer, CPU pinning) and a
project header \code{src/memprof.h} (page-policy allocation, ring generators, counters). Raw
per-repetition rows are reduced by \code{tools/summarize} (median, IQR), and \code{scripts/derive.py}
computes every derived number (plateaus, knee, Little's Law, AMAT fit) into \code{summary.csv}, which the
report reads directly.
\begin{knobs}
  \code{mem_latency}   & single-thread pointer-chase latency: working-set size, access size, node stride, page size, ring generator, timer, counters \\
  \code{mem_bandwidth} & streaming bandwidth with $T$ pinned threads: read/write/copy/triad/rwmix, access-size blocks, byte stride \\
  \code{mem_loaded}    & loaded latency: $T$ load-generator cores (stream, rwmix, or $k$ independent chases) plus a concurrent latency probe \\
  \code{kernel_miss}   & SAXPY loop whose footprint, order and page size vary, with counters \\
  \code{clock_check}   & diagnostic: fixed core clock vs.\ wall time lost by a pinned thread, idle DRAM traffic \\
  Core counters        & \code{perf_event_open(2)} called directly, per thread, user space only \cite{linux_perf} \\
  DRAM counters        & uncore memory-controller events \code{data_reads}/\code{data_writes}, system-wide \\
  Intel MLC            & v3.13, run as a normal user (\code{run.sh mlc}) \cite{intel_mlc} \\
\end{knobs}
We call \code{perf_event_open(2)} directly (not \code{perf}, \sysPerf) so counts cover only each timed
region. Core events (cycles, instructions, L1D read misses, LLC read references and misses, dTLB walks) are
counted in user mode for the measuring thread; multiplexed counts are scaled by
\code{time_enabled/time_running}, with the running fraction recorded. DRAM traffic comes from the uncore
memory-controller PMU (\code{uncore_imc} \code{data_reads}/\code{data_writes}, one count per
\qty{64}{\byte} line), opened system-wide around the same timed windows (\code{perf_event_paranoid}~$\le 0$);
the idle-system background is measured separately (\code{clock_check}) and quoted next to each result.
Intel MLC v3.13 is the external cross-check (\cref{sec:validation}); without root it cannot disable the
prefetchers, so it runs with \code{-e -r}. MLC's MB/s ($10^6$\,B/s) are converted to GiB/s ($2^{30}$\,B/s).

\subsection{Controls}
\begin{itemize}[nosep, leftmargin=*]
  \item \textbf{Pinning and isolation:} \acspinning. SMT: \sysSMT. The machine is otherwise idle; all data
        were collected on 2026-10-21, 20:22--20:39, in one \code{run.sh all} invocation.
  \item \textbf{Frequency:} governor \sysGovernor, turbo \sysTurbo, so the core runs at its fixed
        \qty{2.8}{\giga\hertz} base clock (the TSC rate); cycles are ns~$\times$~\valqty{zq_ghz}{\giga\hertz}.
        User-mode cycles over wall time give a median of \valqty{ws_ctr_ghz}{\giga\hertz}, and the in-process
        add-chain probe reads \val{probe_st_min}--\valqty{probe_st_max}{\giga\hertz} in all \val{probe_st_n}
        single-thread repetitions (\val{probe_ld_min}--\valqty{probe_ld_max}{\giga\hertz} in \code{mem_loaded}).
        The probe is not used for conversion, because an ACPI event storm once disturbed it; that event
        (GPE~0x42) was masked, and a pinned thread now keeps \val{cc_userfrac_min}--\val{cc_userfrac} of the
        wall time (\cref{sec:anomalies}).
  \item \textbf{Pages:} every buffer is a \qty{2}{\mebi\byte}-aligned anonymous mapping whose page size is set
        by \code{madvise} \emph{before first touch} (\code{MADV_HUGEPAGE} or \code{MADV_NOHUGEPAGE}; the latter
        is essential because THP is \sysHugePages{}). Each run records the \code{AnonHugePages} granted.
  \item \textbf{Warm-up, repetitions, timing:} buffers are pre-faulted and traversed once before timing; 8
        repetitions per point (median, IQR bars); each timed interval lasts
        $\geq$\qtyrange{40}{200}{\milli\second}, so timer overhead is negligible.
  \item \textbf{Code version and provenance:} the measured sources are byte-identical to commit
        \code{b00db66} (verified with \code{git diff} and file times; the run started from an uncommitted tree
        on top of \code{29b8630}). Raw data and per-run metadata are listed in \cref{app:repro}.
\end{itemize}
